\documentclass[aspectratio=169]{beamer}
\usepackage[utf8]{inputenc}
\usepackage[T1]{fontenc}
\usepackage[es-ES]{babel}
\usepackage{graphicx}
\usepackage{tikz}
\usepackage{listings}
\usepackage{xcolor}

% Tema
\usetheme{Madrid}
\usecolortheme{default}

% Colores personalizados
\definecolor{unsazul}{RGB}{30, 60, 120}
\definecolor{unsverde}{RGB}{0, 100, 0}
\definecolor{unsrojo}{RGB}{200, 100, 100}

\setbeamercolor{structure}{fg=unsazul}
\setbeamercolor{palette primary}{bg=unsazul,fg=white}
\setbeamercolor{palette secondary}{bg=unsazul!80,fg=white}
\setbeamercolor{palette tertiary}{bg=unsazul!60,fg=white}

% Configuración de listings
\lstset{
    basicstyle=\ttfamily\small,
    keywordstyle=\color{unsazul}\bfseries,
    commentstyle=\color{gray}\itshape,
    stringstyle=\color{unsverde},
    showstringspaces=false,
    breaklines=true,
    frame=single,
    backgroundcolor=\color{gray!10}
}

% Información del documento
\title{Métodos de Análisis de Datos 1}
\subtitle{Clase 4: Calidad de Datos}
\author{Universidad Nacional del Sur}
\date{}

\begin{document}

% Portada
\begin{frame}
\titlepage
\end{frame}

% Contenidos
\begin{frame}{Contenidos de hoy}
\tableofcontents
\end{frame}

\section{Calidad de datos}

\begin{frame}{¿Qué es la calidad de datos?}
\begin{block}{Definición}
Los datos de calidad son aquellos que son \textbf{aptos para su uso previsto} en operaciones, toma de decisiones y planificación.
\end{block}

\pause

\begin{block}{Dimensiones de la calidad}
\begin{enumerate}
    \item \textbf{Completitud:} ¿Hay valores faltantes?
    \item \textbf{Consistencia:} ¿Los datos son coherentes entre sí?
    \item \textbf{Precisión:} ¿Los datos son correctos?
    \item \textbf{Actualidad:} ¿Los datos están actualizados?
    \item \textbf{Validez:} ¿Los datos cumplen reglas de negocio?
\end{enumerate}
\end{block}
\end{frame}

\begin{frame}{La realidad de los datos}
\begin{alertblock}{Dato importante}
En la práctica, \textbf{la mayoría de los datos están sucios}
\end{alertblock}

\begin{itemize}
    \item 60-80\% del tiempo de un proyecto se gasta en limpieza de datos
    \item Los datos crudos casi nunca están listos para análisis
    \item La calidad de los datos determina la calidad del análisis
\end{itemize}

\pause

\vspace{1em}
\begin{exampleblock}{Regla de oro}
\textbf{Garbage in, garbage out}

Si los datos de entrada son malos, los resultados también lo serán
\end{exampleblock}
\end{frame}

\section{Valores faltantes}

\begin{frame}{Valores faltantes: El problema más común}
\begin{block}{¿Qué son?}
Valores que deberían estar presentes pero no lo están
\end{block}

\begin{block}{Representaciones comunes}
\begin{itemize}
    \item Python: \texttt{NaN}, \texttt{None}
    \item R: \texttt{NA}, \texttt{NaN}, \texttt{NULL}
    \item CSV: celdas vacías, \texttt{NA}, \texttt{?}, \texttt{-999}
\end{itemize}
\end{block}

\pause

\begin{alertblock}{¿Por qué faltan?}
\begin{itemize}
    \item Error en la recolección
    \item Sensor/dispositivo falló
    \item Usuario no respondió
    \item No aplica (ej: salario de un desempleado)
    \item Datos perdidos en transferencia
\end{itemize}
\end{alertblock}
\end{frame}

\begin{frame}{Tipos de valores faltantes}
\begin{block}{MCAR: Missing Completely At Random}
La ausencia \textbf{no tiene relación} con ninguna variable

\textit{Ejemplo:} Un sensor falla aleatoriamente
\end{block}

\pause

\begin{block}{MAR: Missing At Random}
La ausencia está relacionada con \textbf{variables observadas}

\textit{Ejemplo:} Personas jóvenes no responden pregunta sobre ingresos (relacionado con edad, no con ingresos)
\end{block}

\pause

\begin{block}{MNAR: Missing Not At Random}
La ausencia está relacionada con el \textbf{valor faltante mismo}

\textit{Ejemplo:} Personas con ingresos altos no responden pregunta sobre ingresos
\end{block}

\vspace{0.5em}
\small{El tipo de missingness determina qué estrategia usar}
\end{frame}

\begin{frame}[fragile]{Detección de valores faltantes}
\begin{columns}[t]
\column{0.5\textwidth}
\textbf{Python:}
\begin{lstlisting}[language=Python]
# Contar NAs por columna
df.isnull().sum()

# Porcentaje
(df.isnull().sum() / len(df)) * 100

# Visualizar patrón
import missingno as msno
msno.matrix(df)
msno.heatmap(df)
\end{lstlisting}

\column{0.5\textwidth}
\textbf{R:}
\begin{lstlisting}[language=R]
# Contar NAs
colSums(is.na(df))

# Porcentaje
colMeans(is.na(df)) * 100

# Visualizar
library(naniar)
vis_miss(df)
gg_miss_var(df)
\end{lstlisting}
\end{columns}
\end{frame}

\begin{frame}{Estrategias para valores faltantes}
\begin{block}{1. Eliminación}
Borrar filas o columnas con valores faltantes
\end{block}

\begin{block}{2. Imputación simple}
Reemplazar con media, mediana, moda o constante
\end{block}

\begin{block}{3. Imputación avanzada}
Usar algoritmos para predecir valores faltantes (KNN, MICE, regresión)
\end{block}

\vspace{1em}
\begin{alertblock}{¿Cuál elegir?}
Depende de:
\begin{itemize}
    \item Porcentaje de valores faltantes
    \item Tipo de missingness (MCAR, MAR, MNAR)
    \item Importancia de la variable
    \item Contexto del problema
\end{itemize}
\end{alertblock}
\end{frame}

\begin{frame}{Estrategia 1: Eliminación}
\begin{columns}[t]
\column{0.5\textwidth}
\textbf{Ventajas:}
\begin{itemize}
    \item Simple
    \item Rápido
    \item No introduce sesgo si es MCAR
\end{itemize}

\column{0.5\textwidth}
\textbf{Desventajas:}
\begin{itemize}
    \item Pérdida de información
    \item Reduce tamaño del dataset
    \item Sesgos si no es MCAR
\end{itemize}
\end{columns}

\vspace{1em}

\begin{alertblock}{Regla general}
Solo eliminar si:
\begin{itemize}
    \item Menos del 5\% de valores faltantes
    \item Los datos son MCAR
    \item El dataset es suficientemente grande
\end{itemize}
\end{alertblock}
\end{frame}

\begin{frame}[fragile]{Eliminación: Código}
\begin{columns}[t]
\column{0.5\textwidth}
\textbf{Python:}
\begin{lstlisting}[language=Python]
# Eliminar filas con cualquier NA
df_clean = df.dropna()

# Eliminar filas donde 
# columnas específicas tienen NA
df_clean = df.dropna(
    subset=['col1', 'col2']
)

# Eliminar columnas 
# con >50% de NAs
umbral = len(df) * 0.5
df_clean = df.dropna(
    thresh=umbral, 
    axis=1
)
\end{lstlisting}

\column{0.5\textwidth}
\textbf{R:}
\begin{lstlisting}[language=R]
# Eliminar filas con NA
df_clean <- na.omit(df)

# Columnas específicas
df_clean <- df %>%
  filter(
    !is.na(col1), 
    !is.na(col2)
  )

# Eliminar columnas 
# con >50% NAs
umbral <- nrow(df) * 0.5
df_clean <- df %>%
  select(where(
    ~ sum(is.na(.)) < umbral
  ))
\end{lstlisting}
\end{columns}
\end{frame}

\begin{frame}{Estrategia 2: Imputación simple}
\begin{block}{Opciones}
\begin{itemize}
    \item \textbf{Media:} Para variables numéricas, distribución simétrica
    \item \textbf{Mediana:} Para variables numéricas con outliers
    \item \textbf{Moda:} Para variables categóricas
    \item \textbf{Constante:} 0, "desconocido", etc.
\end{itemize}
\end{block}

\pause

\begin{columns}[t]
\column{0.5\textwidth}
\textbf{Ventajas:}
\begin{itemize}
    \item Simple y rápido
    \item No reduce tamaño
    \item Funciona razonablemente bien
\end{itemize}

\column{0.5\textwidth}
\textbf{Desventajas:}
\begin{itemize}
    \item Subestima varianza
    \item No considera relaciones
    \item Puede introducir sesgo
\end{itemize}
\end{columns}
\end{frame}

\begin{frame}[fragile]{Imputación simple: Código}
\begin{columns}[t]
\column{0.5\textwidth}
\textbf{Python:}
\begin{lstlisting}[language=Python]
from sklearn.impute import SimpleImputer

# Media
imputer = SimpleImputer(
    strategy='mean'
)
df[['col1', 'col2']] = imputer.fit_transform(
    df[['col1', 'col2']]
)

# Moda (más frecuente)
imputer = SimpleImputer(
    strategy='most_frequent'
)

# Constante
df['col'].fillna(0, inplace=True)
\end{lstlisting}

\column{0.5\textwidth}
\textbf{R:}
\begin{lstlisting}[language=R]
library(tidyverse)

# Media
df <- df %>%
  mutate(col1 = ifelse(
    is.na(col1), 
    mean(col1, na.rm = TRUE), 
    col1
  ))

# Mediana (todas numéricas)
df <- df %>%
  mutate(across(
    where(is.numeric), 
    ~replace_na(., median(., na.rm = TRUE))
  ))
\end{lstlisting}
\end{columns}
\end{frame}

\begin{frame}{Estrategia 3: Imputación avanzada}
\begin{block}{Métodos}
\begin{itemize}
    \item \textbf{KNN:} Imputa usando valores de observaciones similares
    \item \textbf{Regresión:} Predice valores faltantes usando otras variables
    \item \textbf{MICE:} Multiple Imputation by Chained Equations
    \item \textbf{Interpolación:} Para series temporales
\end{itemize}
\end{block}

\pause

\begin{alertblock}{¿Cuándo usar?}
\begin{itemize}
    \item Cuando hay relaciones fuertes entre variables
    \item Dataset mediano/grande
    \item Valores faltantes MAR
    \item Análisis requiere alta precisión
\end{itemize}
\end{alertblock}
\end{frame}

\begin{frame}[fragile]{KNN Imputation}
\begin{block}{Idea}
Usar los K vecinos más cercanos para estimar el valor faltante
\end{block}

\textbf{Python:}
\begin{lstlisting}[language=Python]
from sklearn.impute import KNNImputer

# Imputar con 5 vecinos
imputer = KNNImputer(n_neighbors=5)
df_imputed = pd.DataFrame(
    imputer.fit_transform(df),
    columns=df.columns
)
\end{lstlisting}

\textbf{R:}
\begin{lstlisting}[language=R]
library(mice)

# MICE (más sofisticado que KNN simple)
imputed <- mice(df, m = 5, method = 'pmm', seed = 42)
df_complete <- complete(imputed, 1)
\end{lstlisting}
\end{frame}

\section{Duplicados}

\begin{frame}{Duplicados: Registros repetidos}
\begin{block}{¿Qué son duplicados?}
Filas que aparecen más de una vez en el dataset
\end{block}

\begin{block}{¿Por qué ocurren?}
\begin{itemize}
    \item Error en la recolección (doble entrada)
    \item Múltiples fuentes sin deduplicación
    \item Problemas en joins/merges
    \item Errores de sistema
\end{itemize}
\end{block}

\pause

\begin{alertblock}{Impacto}
\begin{itemize}
    \item Sesgan estadísticas (media, totales)
    \item Inflan el tamaño del dataset
    \item Problemas en modelos de ML
    \item Resultados incorrectos
\end{itemize}
\end{alertblock}
\end{frame}

\begin{frame}[fragile]{Detección de duplicados}
\begin{columns}[t]
\column{0.5\textwidth}
\textbf{Python:}
\begin{lstlisting}[language=Python]
# Todas las columnas
duplicados = df.duplicated()
print(f"Total: {duplicados.sum()}")

# Columnas específicas
duplicados = df.duplicated(
    subset=['id', 'fecha']
)

# Ver duplicados
df[duplicados]
\end{lstlisting}

\column{0.5\textwidth}
\textbf{R:}
\begin{lstlisting}[language=R]
# Detectar duplicados
duplicados <- duplicated(df)
cat("Total:", sum(duplicados))

# Columnas específicas
duplicados <- df %>%
  duplicated(by = c('id', 'fecha'))

# Ver duplicados
df[duplicados, ]
\end{lstlisting}
\end{columns}
\end{frame}

\begin{frame}[fragile]{Eliminación de duplicados}
\begin{columns}[t]
\column{0.5\textwidth}
\textbf{Python:}
\begin{lstlisting}[language=Python]
# Eliminar (mantiene primero)
df_clean = df.drop_duplicates()

# Por columnas específicas
df_clean = df.drop_duplicates(
    subset=['id', 'fecha']
)

# Mantener último
df_clean = df.drop_duplicates(
    keep='last'
)

# Ver cuáles se eliminan
df_clean = df.drop_duplicates(
    keep=False
)  # Elimina todos
\end{lstlisting}

\column{0.5\textwidth}
\textbf{R:}
\begin{lstlisting}[language=R]
# Eliminar duplicados
df_clean <- distinct(df)

# Columnas específicas
df_clean <- df %>%
  distinct(id, fecha, 
           .keep_all = TRUE)
\end{lstlisting}
\end{columns}
\end{frame}

\section{Inconsistencias}

\begin{frame}{Inconsistencias: Datos que no tienen sentido}
\begin{block}{Tipos de inconsistencias}
\begin{enumerate}
    \item \textbf{Valores fuera de rango:} Edad = -5, temperatura = 500°C
    \item \textbf{Formatos incorrectos:} Fecha "32/13/2024"
    \item \textbf{Inconsistencias entre variables:} Fecha nacimiento > Fecha graduación
    \item \textbf{Errores tipográficos:} "Bahia Blanca", "B. Blanca", "Bahía blanca"
    \item \textbf{Valores imposibles:} Porcentaje > 100\%, probabilidad < 0
\end{enumerate}
\end{block}
\end{frame}

\begin{frame}[fragile]{Detección de inconsistencias}
\begin{columns}[t]
\column{0.5\textwidth}
\textbf{Python:}
\begin{lstlisting}[language=Python]
# Valores fuera de rango
print("Edades negativas:", 
      (df['edad'] < 0).sum())
print("Edades >120:", 
      (df['edad'] > 120).sum())

# Rangos
print(f"Min: {df['edad'].min()}")
print(f"Max: {df['edad'].max()}")

# Valores únicos
print(df['ciudad'].value_counts())

# Entre columnas
inconsist = df[
    df['fecha_nac'] > df['fecha_reg']
]
\end{lstlisting}

\column{0.5\textwidth}
\textbf{R:}
\begin{lstlisting}[language=R]
# Fuera de rango
cat("Negativas:", 
    sum(df$edad < 0, na.rm = TRUE))

# Rangos
cat("Min:", min(df$edad, na.rm = TRUE))
cat("Max:", max(df$edad, na.rm = TRUE))

# Valores únicos
table(df$ciudad)

# Entre columnas
inconsist <- df %>%
  filter(fecha_nac > fecha_reg)
\end{lstlisting}
\end{columns}
\end{frame}

\begin{frame}[fragile]{Corrección: Estandarización de texto}
\begin{columns}[t]
\column{0.5\textwidth}
\textbf{Python:}
\begin{lstlisting}[language=Python]
# Lowercase y trim
df['ciudad'] = df['ciudad'].str.lower().str.strip()

# Reemplazar variaciones
df['ciudad'] = df['ciudad'].replace({
    'b. blanca': 'bahia blanca',
    'bblanca': 'bahia blanca'
})

# Regex para reemplazar
df['ciudad'] = df['ciudad'].str.replace(
    r'bahia\s+blanca', 
    'bahia blanca',
    regex=True
)
\end{lstlisting}

\column{0.5\textwidth}
\textbf{R:}
\begin{lstlisting}[language=R]
# Lowercase y trim
df <- df %>%
  mutate(ciudad = str_to_lower(
    str_trim(ciudad)
  ))

# Reemplazar
df <- df %>%
  mutate(ciudad = recode(ciudad,
    'b. blanca' = 'bahia blanca',
    'bblanca' = 'bahia blanca'
  ))
\end{lstlisting}
\end{columns}
\end{frame}

\begin{frame}[fragile]{Corrección: Valores fuera de rango}
\begin{columns}[t]
\column{0.5\textwidth}
\textbf{Python:}
\begin{lstlisting}[language=Python]
# Reemplazar con NA
df.loc[df['edad'] < 0, 'edad'] = np.nan
df.loc[df['edad'] > 120, 'edad'] = np.nan

# Clip (limitar valores)
df['edad'] = df['edad'].clip(
    lower=0, 
    upper=120
)

# Filtrar filas válidas
df_clean = df[
    (df['edad'] >= 0) & 
    (df['edad'] <= 120)
]
\end{lstlisting}

\column{0.5\textwidth}
\textbf{R:}
\begin{lstlisting}[language=R]
# Reemplazar con NA
df <- df %>%
  mutate(edad = ifelse(
    edad < 0 | edad > 120, 
    NA, 
    edad
  ))

# Con pmin/pmax
df <- df %>%
  mutate(edad = pmin(
    pmax(edad, 0), 
    120
  ))
\end{lstlisting}
\end{columns}
\end{frame}

\section{Resumen}

\begin{frame}{Pipeline de limpieza de datos}
\begin{enumerate}
    \item \textbf{Cargar datos}
    \item \textbf{Exploración inicial:} dimensiones, tipos, primeras filas
    \item \textbf{Detectar valores faltantes:} porcentaje, patrón
    \item \textbf{Decidir estrategia:} eliminar o imputar
    \item \textbf{Detectar duplicados:} por todas las columnas o subset
    \item \textbf{Eliminar duplicados}
    \item \textbf{Detectar inconsistencias:} rangos, formatos, relaciones
    \item \textbf{Corregir inconsistencias:} estandarizar, clip, reemplazar
    \item \textbf{Verificación final:} re-explorar datos limpios
    \item \textbf{Documentar:} qué se hizo y por qué
\end{enumerate}
\end{frame}

\begin{frame}{Resumen de la clase}
\begin{itemize}
    \item La \textbf{calidad de datos} tiene múltiples dimensiones
    \item Los \textbf{valores faltantes} son el problema más común
    \item Tipos de missingness: \textbf{MCAR, MAR, MNAR}
    \item Estrategias: \textbf{eliminación} (simple) o \textbf{imputación} (simple/avanzada)
    \item Los \textbf{duplicados} sesgan resultados y deben eliminarse
    \item Las \textbf{inconsistencias} requieren análisis cuidadoso
    \item La \textbf{estandarización} de texto es crucial
    \item \textbf{Documentar} todas las decisiones de limpieza
\end{itemize}

\vspace{1em}

\begin{block}{Próxima clase}
\textbf{Transformación de datos:} tipos, variables derivadas, codificación de categóricas
\end{block}
\end{frame}

\begin{frame}{Tarea}
\begin{enumerate}
    \item Tomar el dataset de la clase anterior
    \item Realizar análisis de calidad:
    \begin{itemize}
        \item ¿Hay valores faltantes? ¿Cuántos? ¿En qué columnas?
        \item ¿Hay duplicados?
        \item ¿Hay valores fuera de rango?
        \item ¿Hay inconsistencias en texto?
    \end{itemize}
    \item Aplicar estrategias de limpieza apropiadas
    \item Documentar en el notebook:
    \begin{itemize}
        \item Qué problemas encontraste
        \item Qué estrategias usaste
        \item Por qué elegiste esas estrategias
        \item Comparar antes vs después
    \end{itemize}
    \item Subir a GitHub
\end{enumerate}
\end{frame}

\begin{frame}[standout]
\Huge ¿Preguntas?
\end{frame}

\end{document}
